% C9: Dynamic Model Onboarding — Controlled Family Sharing Experiment
% Cross-references: D.3 (full figure), C3 (cold-start ablation), C6 (learning curves)

\subsection{Dynamic Model Onboarding}
\label{appendix:model_onboarding}

Production LLM portfolios evolve as providers release new models.
This section evaluates \texttt{register\_model()}'s ability to integrate a
newcomer into a running router, using a \textbf{controlled experiment}
that isolates the effect of hybrid family sharing from all other
onboarding mechanisms.

\subsubsection{Experimental Design}

The same newcomer model (\texttt{openai/gpt-4.1}) is introduced under
two family-assignment conditions:

\begin{itemize}[nosep]
    \item \textbf{Shared family:} \texttt{gpt-4.1} joins the existing
          \texttt{openai/gpt-4} family alongside \texttt{gpt-4-turbo}
          (auto-inferred via \texttt{infer\_model\_family()}).
          The newcomer inherits the family's shared $\betavec$ via
          the Hybrid LinUCB architecture.
    \item \textbf{Isolated family:} \texttt{gpt-4.1} is forced into a
          singleton family (\texttt{openai/gpt-4.1-isolated}) via
          \texttt{family\_map} override.  No family sharing occurs.
\end{itemize}

\noindent Crucially, \emph{all other variables are held constant}:
same model, same data, same seed schedule, same embeddings, and---most
importantly---the same semantic bootstrapping
(\texttt{gpt-4.1} $\to$ \texttt{gpt-4-turbo}, cosine similarity $0.849$,
$N_{\text{eff}}{=}5$).  The only treatment variable is whether the
newcomer's family $\betavec$ is shared with \texttt{gpt-4-turbo}.

Both newcomers are introduced at step~501 into a router that has
completed 500~steps of burn-in on the 2-model base fleet
(\texttt{gpt-4-turbo} + \texttt{mixtral-8x7b}), using production
warmup priors.  Routing and learning continue through step~1,121.
Holdout evaluation ($N{=}750$) uses the final router state.
All results: 20 seeds (42--61), 95\% CIs ($t_{19}$).

\subsubsection{Reference Baselines}

\begin{table}[h]
\centering
\caption{Static reference points on the holdout set (3-model pool).}
\label{tab:onboarding_baselines}
\small
\begin{tabular}{@{}lc@{}}
\toprule
\textbf{Baseline} & \textbf{Holdout Reward} \\
\midrule
Oracle (per-prompt best of 3) & 1.000 \\
Always-newcomer (gpt-4.1 only) & 0.983 \\
\bottomrule
\end{tabular}
\end{table}

\subsubsection{Holdout Quality}

\begin{table}[h]
\centering
\caption{Model onboarding results ($\lambda{=}0$, 20 seeds, 95\% CI).
Family sharing provides no significant holdout quality advantage when
semantic bootstrapping is strong.}
\label{tab:model_onboarding}
\small
\begin{tabular}{@{}lccc@{}}
\toprule
\textbf{Condition} & \textbf{Holdout Reward} & \textbf{95\% CI}
  & \textbf{Newcomer Routed} \\
\midrule
Shared family   & 0.965 & $\pm$0.005 & 448 $\pm$ 33 \\
Isolated family & 0.965 & $\pm$0.003 & 483 $\pm$ 19 \\
\midrule
\textbf{Difference} & +0.000 & & $-$35 \\
\bottomrule
\end{tabular}
\end{table}

\noindent A paired $t$-test on per-trial holdout rewards confirms no
significant difference: $t(19){=}0.05$, $p{=}0.958$, with 12 of 20
seeds favoring the shared condition and 8 favoring isolated.

\subsubsection{Prediction Error Convergence}

While final quality is indistinguishable, calibration \emph{speed}
differs.  We track the corralling-weighted prediction error (MAE between
predicted and actual normalized reward) for prompts routed to the
newcomer.

\begin{table}[h]
\centering
\caption{Newcomer prediction error convergence (first and last 50
selections, averaged over 10 bins of 5).}
\label{tab:pred_error_convergence}
\small
\begin{tabular}{@{}lccc@{}}
\toprule
\textbf{Condition} & \textbf{First 50 MAE} & \textbf{Last 50 MAE}
  & \textbf{Reduction} \\
\midrule
Shared family   & 0.090 & 0.046 & 49\% \\
Isolated family & 0.165 & 0.045 & 73\% \\
\bottomrule
\end{tabular}
\end{table}

The shared-family condition starts with 46\% lower initial prediction
error (0.090 vs.\ 0.165), confirming that the inherited $\betavec$
provides a better initial estimate.  However, both conditions converge
to nearly identical final MAE (${\sim}0.045$), explaining the null
result on holdout quality.

\subsubsection{Key Findings}

\begin{enumerate}[nosep]
    \item \textbf{Family sharing does not improve final quality.}
          When semantic bootstrapping is strong (similarity $0.849$),
          the additional family $\betavec$ provides no statistically
          significant holdout quality advantage ($p{=}0.958$).
          The per-model $\thetavec$ initialized from the semantic
          neighbor captures sufficient signal.
    \item \textbf{Family sharing accelerates early calibration.}
          The shared $\betavec$ halves the initial prediction error
          (0.090 vs.\ 0.165), providing a better warm start for the
          first ${\sim}50$ routing decisions.  This advantage dissipates
          as the per-model $\thetavec$ accumulates observations.
    \item \textbf{Semantic bootstrapping is the dominant mechanism.}
          Both conditions receive identical semantic transfer from
          \texttt{gpt-4-turbo} ($\thetavec$ initialization,
          $N_{\text{eff}}{=}5$).  This is sufficient for the router to
          quickly calibrate the newcomer, regardless of family assignment.
    \item \textbf{Both conditions succeed comprehensively.}
          Holdout reward ($0.965$) substantially exceeds the 2-model
          baseline ($0.911$, Table~\ref{tab:learning_curve}) and
          approaches the always-newcomer ceiling ($0.983$),
          confirming that \texttt{register\_model()} successfully
          integrates a new model regardless of family configuration.
\end{enumerate}

\subsubsection{Practical Implications}

\paragraph{Just call \texttt{register\_model()}.}
The operational message for practitioners is simple:
\texttt{register\_model()} is sufficient for effective model onboarding.
The API automatically discovers the nearest semantic neighbor, bootstraps
the newcomer's reward estimates, and resets confidence to ensure
exploration---no manual family configuration required.
Adding a third model to a 2-model fleet improved holdout quality from
$0.911$ to $0.965$ (+5.4~pp), and this improvement was identical
regardless of whether family sharing was active ($p{=}0.958$).
The newcomer reaches full calibration within ${\sim}600$ routing
decisions; at typical production throughput ($>$100~req/hr), this
completes within a few hours.

\paragraph{When family sharing matters.}
The hybrid architecture's family $\betavec$ mechanism provides faster
early calibration (46\% lower initial prediction error for the first
${\sim}50$ decisions), which is valuable in latency-sensitive
deployments where even the first 50 routing decisions matter---for
example, a low-traffic service where 50 requests may take hours.
For most production workloads with reasonable throughput, the early
calibration advantage dissipates before it has operational impact.
Family sharing becomes more relevant when: (a)~no strong semantic
neighbor exists (e.g., a model from a new provider with no similar
models in the fleet), (b)~the family has accumulated extensive routing
history that a newcomer can leverage, or (c)~the model portfolio is
large enough ($K \geq 8$) that per-model learning is slow due to
exploration budget spreading across many arms.

\paragraph{Contrast with the semantic transfer null result.}
Section~\ref{sec:transfer_null} reports a null result for \emph{per-model}
semantic $\thetavec$ transfer in a leave-one-out design.  That experiment
evaluated whether transferring $\thetavec$ from a neighbor provides
meaningful advantage \emph{when the newcomer is a permanent member of
the original pool}---and found it does not, because established models'
accumulated precision matrices ($\Amat$) dominate UCB scores, starving
the newcomer of traffic.  The current experiment evaluates \emph{dynamic}
onboarding with \texttt{register\_model()}, where semantic bootstrapping
\emph{does} provide effective initialization.  The distinction is that
\texttt{register\_model()} combines semantic transfer with confidence
reset ($\Amat{=}\mathbf{I}$), ensuring the newcomer explores rather
than being starved.  This confidence reset is the critical ingredient:
it gives the newcomer a comparable UCB bonus to established models,
allowing the router to learn the newcomer's true quality within a few
hundred observations.
